\section{Life as Art}

\begin{quotex}
A human being should be able to change a diaper, plan an invasion, butcher a hog, con a ship, design a building, write a sonnet, balance accounts, build a wall, set a bone, comfort the dying, take orders, give orders, cooperate, act alone, solve equations, analyze a new problem, pitch manure, program a computer, cook a tasty meal, fight efficiently, die gallantly. Specialization is for insects.
\flright{\textsc{Robert Heinlein}}
\end{quotex}


\begin{quotex}
Creation emanates from the potentialities inherent in God's being, a being which radiates ``into'' the coming into visibility and intelligibility of all that surrounds us.
\flright{\textsc{George Steiner}, \emph{Grammars of Creation}}
\end{quotex}


\paragraph{Depth and Breadth of Thought}


\begin{quotationx}
Thinking is never a sharp, neat, linear process; it could rather be compared to the progress of a boat on a lake.

When you day-dream you drift before the wind; when you read or listen to a narrative you travel like a barge towed by a tug. But in each case the progress of the boat causes ripples on the lake spreading in all directions --- memories, images, associations; some of these move quicker than the boat itself and create anticipation; others penetrate into the deep. The boat symbolized focal awareness, the ripples on the surface are the fringes of consciousness, and you can furnish the deeps, according to taste. Which the nasty eddies of repressed complexes, the deep-water currents of the collective unconscious, or with archetypal coral-reefs. When thinking is in the tow of a narrative, focal awareness must stick to its course and cannot follow the ripples on their journey across the lake; but is their presence all round the horizon, on the peripheries of awareness, which provides resonance, coulour, and depth, the atmosphere and feel of the story.
\flright{\textsc{Arthur Koestler}, \emph{The Act of Creation}}
\end{quotationx}


\paragraph{Life as Art}


Non-duality is not a set of doctrines to feed the head; nor is it mere sentimentality to satisfy the heart. It is rather an understanding that can only be lived, not just thought about. We may be tempted to idolize the scientist or artist because their experiences are so compelling and the results so transforming, yet creativity is not just for the artist or the scientist, but it a way of life for each one of us. One's own life can be lived as a work of art ... consider this suggestion from Nietzsche:

\begin{quotationx}
\textbf{One thing is needful.} --- To ``give style'' to one's character--- a great and rare art! It is practiced by those who survey all the strengths and weaknesses of their nature and then fit them into an artistic plan until every one of them appears as art and reason and even weaknesses delight the eye. Here a large mass of second nature has been added; there a piece of original nature has been removed --- both times through long practice and daily work at it. Here the ugly that could not be removed is concealed; there it has been reinterpreted and made sublime. Much that is vague and resisted shaping has been saved and exploited for distant views; it is meant to beckon toward the far and immeasurable. In the end, when the work is finished, it becomes evident how the constraint of a single taste governed and formed everything large and small. Whether this taste was good or bad is less important than one might suppose, if only it was a single taste!

It will be the strong and domineering natures that enjoy their finest gaiety in such constraint and perfection under a law of their own; the passion of their tremendous will relaxes in the face of all stylized nature, of all conquered and serving nature. Even when they have to build palaces and design gardens they demur at giving nature freedom.

Conversely, it is the weak characters without power over themselves that hate the constraint of style. They feel that if this bitter and evil constraint were imposed upon them they would be demeaned; they become slaves as soon as they serve; they hate to serve. Such spirits --- and they may be of the first rank --- are always out to shape and interpret their environment as free nature: wild, arbitrary, fantastic, disorderly, and surprising. And they are well advised because it is only in this way that they can give pleasure to themselves. For one thing is needful: that a human being should attain satisfaction with himself, whether it be by means of this or that poetry or art; only then is a human being at all tolerable to behold. Whoever is dissatisfied with himself is continually ready for revenge, and we others will be his victims, if only by having to endure his ugly sight. For the sight of what is ugly makes one bad and gloomy.
\end{quotationx}



\flrightit{Posted on 2006-10-29 by Cologero}

\begin{center}* * *\end{center}

\begin{footnotesize}
\begin{sffamily}
\texttt{Lyon on 2018-05-06 at 16:39 said:}

``Whoever is dissatisfied with himself is continually ready for revenge, and we others will be his victims, if only by having to endure his ugly sight.'' Friedrich Nietzsche

Precisely!


\hfill

\end{sffamily}
\end{footnotesize}

